\documentclass[11pt]{amsart}

\usepackage[margin=1.1in]{geometry}
\usepackage{amsmath,amssymb,amsthm,mathtools}
\usepackage{enumitem}
\usepackage[hidelinks]{hyperref}
\usepackage[capitalize,nameinlink]{cleveref}

\numberwithin{equation}{section}

\theoremstyle{plain}
\newtheorem{theorem}{Theorem}[section]
\newtheorem{proposition}[theorem]{Proposition}
\newtheorem{lemma}[theorem]{Lemma}
\newtheorem{corollary}[theorem]{Corollary}
\newtheorem{fact}{Fact}
\renewcommand{\thefact}{F\arabic{fact}}
\theoremstyle{definition}
\newtheorem{definition}[theorem]{Definition}
\newtheorem{remark}[theorem]{Remark}
\crefname{fact}{Fact}{Facts}

\newcommand{\Z}{\mathbb{Z}}
\newcommand{\R}{\mathbb{R}}
\newcommand{\C}{\mathbb{C}}
\newcommand{\HH}{\mathbb{H}}
\newcommand{\poch}[3]{(#1;#2)_{#3}}
\newcommand{\cert}{\mathcal{W}}
\newcommand{\muh}{\widehat{\mu}}
\newcommand{\Sh}{\widehat{S}}
\newcommand{\fh}{\widehat{f}}
\newcommand{\SXXI}{S_{\mathrm{XXI}}}
\newcommand{\sgn}{\operatorname{sgn}}
\newcommand{\Img}{\operatorname{Im}}

\begin{document}

\title{Evaluation of Au's Watson series and its mixed mock modular completion}
\author{Sumit Kumar Jha}
\address{Greater Noida, India}
\email{sumitkumarjha75@gmail.com}
\subjclass[2020]{Primary 33D15; Secondary 11F37, 11F03, 33F10}
\keywords{Basic hypergeometric series, Wilf--Zeilberger pairs, Watson's
transformation, Ramanujan's series for $1/\pi$, mock modular forms, Zwegers'
$\mu$-function}

\begin{abstract}
Ramanujan's series $\sum_{n\ge0}\frac{(1/2)_n^3}{n!^3\,64^n}(42n+5)=\frac{16}{\pi}$
has many $q$-analogues. Au constructed one of them from a Watson-type
Wilf--Zeilberger seed and asked whether it is mock modular. We evaluate it in
closed form: it equals an infinite product times the bilateral Lambert series
$\sum_{r\in\Z}q^{4r(r+1)}/(1-q^{4r+1})$. The proof has three steps. A
Wilf--Zeilberger pair, whose certificate we give explicitly, turns each partial
sum into a finite sum. Watson's transformation rewrites that finite sum with
nonnegative terms. Re-centering the index and letting the length grow then
produces a two-sided series, which a classical identity converts into the
Lambert series. The Lambert series is a value of Zwegers' $\mu$-function, so
the series is an eta quotient times a mock theta function of weight $1/2$. We
write down its completion, which transforms with weight $-1$ and character
$\chi_{-4}$ on $\Gamma_0(4)$, and we compute the shadow of the weight-$1/2$
factor, $-\frac12\eta(2\tau)^3$. We also prove directly that the $q$-series
recovers $16/\pi$ as $q\to1$, and compare the summand with Campbell's
$q$-analogue of the same formula.
\end{abstract}

\maketitle

\section{Introduction}\label{sec:intro}

In 1914 Ramanujan listed seventeen rapidly converging series for $1/\pi$.
One of them is
\begin{equation}\label{eq:classical}
\sum_{n=0}^{\infty}\frac{(1/2)_n^3}{(n!)^3\,64^n}\,(42n+5)=\frac{16}{\pi},
\end{equation}
where $(a)_0=1$ and $(a)_n=a(a+1)\cdots(a+n-1)$ for $n\ge1$; this is recorded
as \cite[(1.3)]{Au}. A \emph{$q$-analogue} of such a formula is a series in an
extra variable $q$ whose terms tend, after a suitable normalization, to the
terms of the classical series as $q\to1$, and which can itself be summed in
closed form.

Au \cite{Au} found a systematic way to produce $q$-analogues. He starts from a
\emph{Wilf--Zeilberger seed}, a hypergeometric expression from which a whole
family of Wilf--Zeilberger pairs (see \cref{sec:wz}) can be generated. Many of
his $q$-analogues evaluate to infinite products, which are modular forms. Some
evaluate instead to \emph{mock} modular forms, functions that are modular only
after a non-holomorphic correction term is added (see \cref{sec:modular}).
After his Example~XXI, Au writes down a further $q$-analogue of
\eqref{eq:classical}, built from a seed related to Watson's transformation,
and asks whether it too is mock modular \cite[remark after Example~XXI]{Au}.
This paper answers that question. We evaluate the series in closed form,
identify its mock modular part with the one already present in Example~XXI,
give the completion explicitly, and verify the classical limit.

\subsection{The series}\label{sec:series}
For a complex number $a$, a base $r$ with $0<|r|<1$ and an integer $n\ge0$,
put
\[
\poch{a}{r}{n}=\prod_{j=0}^{n-1}(1-ar^j),\qquad
\poch{a}{r}{\infty}=\prod_{j=0}^{\infty}(1-ar^j),
\]
with the empty product equal to $1$. For several parameters we write
$\poch{a_1,\dots,a_s}{r}{n}=\prod_{j}\poch{a_j}{r}{n}$. Throughout, $q$ is a
complex number with $0<|q|<1$, and we use the abbreviations
\[
Q=q^2,\qquad J_d=\poch{q^d}{q^d}{\infty}\quad(d=1,2,3,\dots).
\]
Let $P(q,x)$ be the polynomial
\begin{equation}\label{eq:P}
\begin{aligned}
P(q,x)={}&1+qx+(1+q^2)x^2+2q^3x^3+(3q^2+2q^4)x^4+q^3x^5\\
&+(-2q^2+q^4+q^6)x^6+(-q-q^3+2q^5)x^7-q^4x^8\\
&-(2q^3+3q^5)x^9-2q^4x^{10}-(q^5+q^7)x^{11}-q^6x^{12}-q^7x^{13}.
\end{aligned}
\end{equation}
For each integer $n\ge0$ set
\begin{equation}\label{eq:Bn}
B_n=\frac{\poch{q}{Q}{n}^4\,\poch{Q}{Q^2}{n}^2\,\poch{-Q}{Q}{4n}}
{\poch{Q^2}{Q^2}{2n}^3},\qquad
a_n=\frac{q^{2n}\,P(q,q^{2n})\,B_n}
{(1+q^{8n})(1+q^{2n+1})(1+q^{4n+2})^3},
\end{equation}
and let
\[
S(q)=\sum_{n=0}^{\infty}a_n .
\]
Separating even and odd factors (\cref{fact:evenodd}) shows that $B_n$ can also be
written as
\[
B_n=\frac{\poch{q}{q^2}{n}^4\,\poch{q^2}{q^2}{2n}\,\poch{q^2}{q^4}{n}\,
\poch{q^4}{q^4}{4n}}
{\poch{q^2}{q^2}{4n}\,\poch{q^4}{q^4}{n}\,\poch{q^4}{q^4}{2n}^3},
\]
which is the form in which Au writes the summand. The first term is
\begin{equation}\label{eq:a0}
a_0=\frac{1-q^5}{(1+q)(1+q^2)^2}.
\end{equation}

\subsection{Results}
The main result evaluates $S(q)$.

\begin{theorem}\label{thm:main}
For every complex $q$ with $0<|q|<1$ the series $S(q)$ converges absolutely,
and
\begin{equation}\label{eq:main}
S(q)=\frac{\poch{q}{q^2}{\infty}^2}{J_4^4}\,L(q),\qquad
L(q)=\sum_{r\in\Z}\frac{q^{4r(r+1)}}{1-q^{4r+1}}.
\end{equation}
Both sides extend holomorphically to $q=0$, where they equal $1$.
\end{theorem}

The bilateral series $L(q)$ is the source of mock modularity. Since $L$ also
appears in Au's evaluation of Example~XXI, the two $q$-series are related by an
infinite product (\cref{cor:XXI}).

The next result confirms that $S$ is a $q$-analogue of \eqref{eq:classical};
its proof does not use \cref{thm:main}.

\begin{proposition}\label{prop:classical}
As $q\to1$ through real values $0<q<1$,
\[
\frac{8\,S(q)}{1-q}\longrightarrow\frac{16}{\pi}.
\]
\end{proposition}

Finally, write $q=e^{2\pi i\tau}$ with $\tau$ in the upper half-plane $\HH$.
In \cref{sec:completion} we define an eta quotient $A(\tau)$ of weight $-3/2$,
a function $f(\tau)$ built from Zwegers' $\mu$-function, its completion
$\fh(\tau)=f(\tau)+\frac14R_0(\tau)$ (where $R_0$ is an explicit
non-holomorphic function), and the product $\Sh=A\fh$.

\begin{theorem}\label{thm:completion}
For every $\tau\in\HH$ we have $S(e^{2\pi i\tau})=A(\tau)f(\tau)$. The completed
function $\Sh=A\fh$ satisfies
\[
\Sh\Big(\frac{a\tau+b}{c\tau+d}\Big)=\chi_{-4}(d)\,(c\tau+d)^{-1}\,\Sh(\tau)
\qquad\text{for every }\begin{pmatrix}a&b\\c&d\end{pmatrix}\in\Gamma_0(4),
\]
and the weight-$1/2$ factor $\fh$ is a harmonic Maass form with shadow
\[
\xi_{1/2}\fh=-\tfrac12\,\eta(2\tau)^3 .
\]
In particular $S$ is a mixed mock modular form of weight $-1$.
\end{theorem}

All terms used here ($\Gamma_0(4)$, $\chi_{-4}$, $\xi_k$, harmonic Maass form,
shadow, mixed mock modular form) are defined in \cref{sec:modular}.

\subsection{How the proof works}
The evaluation proceeds in five steps, each short once the right object is
written down.
\begin{enumerate}[label=(\arabic*)]
\item \emph{A finite identity} (\cref{lem:telescope}). We give a two-variable
term $F(n,k)$ and an explicit rational function, the \emph{certificate},
that together form a Wilf--Zeilberger pair. Summing the pair over $k$ shows
that the partial sum $a_0+\dots+a_{N-1}$ equals minus a finite sum $T_N$ of
$2N$ terms.
\item \emph{Watson's transformation} (\cref{lem:watson}). The finite sum $T_N$ is a
terminating very-well-poised series, and Watson's transformation rewrites it as
a sum of $2N$ terms $h_N(k)$ that are all nonnegative when $0<q<1$.
\item \emph{A centred limit} (\cref{lem:centre}). Writing $k=N+r$ and letting
$N\to\infty$, each $h_N(N+r)$ tends to a constant multiple of an explicit term
$b_r$. A bound that does not depend on $N$ justifies exchanging limit and sum.
The result expresses $S(q)$ through the two-sided series $\sum_{r\in\Z}b_r$.
\item \emph{A classical identity} turns $\sum_r b_r$ into the Lambert series
$L(q)$, which proves \cref{thm:main} for $0<q<1$; analytic continuation
gives all $|q|<1$.
\item \emph{Zwegers' theory.} The quotient of $L(q)$ by a theta function is a
special value of Zwegers' $\mu$-function, whose completion and transformation
laws are known. This gives \cref{thm:completion}.
\end{enumerate}
\Cref{prop:classical} is proved separately in \cref{sec:classical} by taking
termwise limits under a uniform bound. \Cref{sec:campbell} compares the
summand with the one in Campbell's $q$-analogue of \eqref{eq:classical}.

\subsection{Computer verification}
Every identity in this paper was checked by computer, using exact rational
arithmetic wherever possible: \cref{thm:main} and \cref{cor:XXI} as
power-series identities modulo $q^{220}$; the certificate identity of
\cref{app:cert} by full polynomial expansion; \cref{lem:telescope,lem:watson,lem:centre}
at many rational values of $q$; and the transformation law and shadow of
\cref{thm:completion} numerically. The scripts are available at
\url{https://github.com/ramanujamath/au-watson-mock-series}.

\section{Notation and background}\label{sec:background}

We write $\Z$, $\R$, $\C$ for the integers, real numbers and complex numbers.
The notation $f=O(g)$ near a stated limit means $|f|\le C|g|$ for some
constant $C$ throughout a neighbourhood of that limit; $O_t$ means that $C$
may depend on $t$.

\subsection{Products and series in \texorpdfstring{$q$}{q}}
The products $\poch{a}{r}{n}$ were defined in \cref{sec:series}. For a positive
integer $j$ we also use the \emph{negative-index} convention
\begin{equation}\label{eq:negindex}
\poch{a}{r}{-j}=\frac{1}{\poch{ar^{-j}}{r}{j}},
\end{equation}
whenever the denominator is nonzero. It is chosen so that
$\poch{a}{r}{m+n}=\poch{a}{r}{m}\poch{ar^m}{r}{n}$ holds for all integers
$m,n$.

\begin{fact}[even and odd factors]\label{fact:evenodd}
For integers $m\ge0$,
\[
\poch{a}{r}{m}\,\poch{-a}{r}{m}=\poch{a^2}{r^2}{m},\qquad
\poch{a}{r}{2m}=\poch{a}{r^2}{m}\,\poch{ar}{r^2}{m}.
\]
The first holds factor by factor since $(1-ar^j)(1+ar^j)=1-a^2r^{2j}$; the
second splits $j$ into even and odd values.
\end{fact}

For $s\ge0$, parameters $a_0,\dots,a_s$, $b_1,\dots,b_s$ and argument $z$, the
\emph{basic hypergeometric series} is
\[
{}_{s+1}\phi_s\!\left[\begin{matrix}a_0,\dots,a_s\\ b_1,\dots,b_s\end{matrix};Q,z\right]
=\sum_{k\ge0}\frac{\poch{a_0,\dots,a_s}{Q}{k}}{\poch{Q,b_1,\dots,b_s}{Q}{k}}\,z^k,
\]
whenever the series terminates or converges and no denominator vanishes. If
one numerator parameter equals $Q^{-m}$ for an integer $m\ge0$, every term
with $k>m$ vanishes and the series \emph{terminates}.

\begin{fact}[Watson's transformation {\cite[(17.9.15)]{DLMF}}]\label{fact:watson}
For an integer $m\ge0$ and any choice of $\sqrt a$,
\begin{multline*}
{}_8\phi_7\!\left[\begin{matrix}a,\ Q\sqrt a,\ -Q\sqrt a,\ b,\ c,\ d,\ e,\ Q^{-m}\\
\sqrt a,\ -\sqrt a,\ aQ/b,\ aQ/c,\ aQ/d,\ aQ/e,\ aQ^{m+1}\end{matrix};
Q,\frac{a^2Q^{m+2}}{bcde}\right]\\
=\frac{\poch{aQ,\ aQ/(de)}{Q}{m}}{\poch{aQ/d,\ aQ/e}{Q}{m}}\;
{}_4\phi_3\!\left[\begin{matrix}aQ/(bc),\ d,\ e,\ Q^{-m}\\
aQ/b,\ aQ/c,\ deQ^{-m}/a\end{matrix};Q,Q\right].
\end{multline*}
(The two parameters $\pm Q\sqrt a$ and $\pm\sqrt a$ occur in pairs, so the
left side does not depend on the sign of $\sqrt a$.)
\end{fact}

\subsection{Wilf--Zeilberger pairs}\label{sec:wz}
Two functions $F(n,k)$ and $G(n,k)$ of integers $n,k$ form a
\emph{Wilf--Zeilberger pair} if
\begin{equation}\label{eq:wzdef}
F(n+1,k)-F(n,k)=G(n,k+1)-G(n,k).
\end{equation}
Summing \eqref{eq:wzdef} over a range of $k$ makes the right side telescope,
so that $\sum_kF(n+1,k)-\sum_kF(n,k)$ is a difference of two boundary values of
$G$. In practice $G=\cert\cdot F$ for a rational function $\cert$ of $q$,
$q^{2n}$ and $q^{2k}$, called the \emph{certificate}. Dividing
\eqref{eq:wzdef} by $F(n,k)$ turns it into an identity between rational
functions, which can be checked by expanding polynomials; this is done in
\cref{app:cert}.

\subsection{Two classical identities}
Put
\[
U(q)=\sum_{r\in\Z}\frac{(-1)^rq^{r^2}}{1-q^{2r-1}},\qquad
\Theta(q)=\sum_{r\in\Z}(-1)^rq^{r^2}.
\]

\begin{fact}[a bilateral identity]\label{fact:LO}
For $0<q<1$,
\[
\sum_{j\ge1}q^{2j}\poch{q}{Q}{j}^2+\sum_{j\ge0}\frac{q^{2j(j-1)}}{\poch{q}{Q}{j}^2}
=-\frac{U(q)}{q\,J_2}.
\]
This is the three-term identity reproduced by Lovejoy and Osburn
\cite[(3.11)]{LO}, with their base replaced by $Q$ and both of their parameters
set equal to $-q$; in this specialization their product prefactor becomes
$1/J_2$, and shifting their bilateral index by one produces the factor
$-q^{-1}$.
\end{fact}

\begin{fact}[Jacobi]\label{fact:jacobi}
$\Theta(q)=J_1^2/J_2$ (a case of the triple product identity), and
\[
\sum_{m\ge0}(-1)^m(2m+1)q^{m(m+1)}=J_2^3 .
\]
\end{fact}

\subsection{Modular forms, mock modular forms and shadows}\label{sec:modular}
Let $\HH=\{\tau\in\C:\Img\tau>0\}$. We write $\tau=u+iy$ with $u,y$ real and
$y>0$, and $q=e^{2\pi i\tau}$; every fractional power $q^\alpha$ means
$e^{2\pi i\alpha\tau}$. The \emph{Dedekind eta function} is
$\eta(\tau)=q^{1/24}\prod_{j\ge1}(1-q^j)$; it has no zeros in $\HH$.

The group $\mathrm{SL}_2(\Z)$ of integer $2\times2$ matrices of determinant
$1$ acts on $\HH$ by $\gamma\tau=(a\tau+b)/(c\tau+d)$ for
$\gamma=\left(\begin{smallmatrix}a&b\\c&d\end{smallmatrix}\right)$. For a
positive integer $M$,
$\Gamma_0(M)=\{\gamma\in\mathrm{SL}_2(\Z):M\mid c\}$. The character
$\chi_{-4}$ is defined by $\chi_{-4}(d)=0$ for even $d$ and
$\chi_{-4}(d)=(-1)^{(d-1)/2}$ for odd $d$. A \emph{cusp} of $\Gamma_0(M)$ is
an orbit of $\Gamma_0(M)$ on $\mathbb{Q}\cup\{\infty\}$; the behaviour of a
function ``at a cusp'' is read off after moving the cusp to $\infty$ by an
element of $\mathrm{SL}_2(\Z)$ and multiplying by the appropriate factor
$(c\tau+d)^{-k}$.

For a weight $k\in\R$ we use the operators
\[
\partial_{\bar\tau}=\tfrac12(\partial_u+i\partial_y),\qquad
\xi_kh=2iy^k\,\overline{\partial_{\bar\tau}h},\qquad
\Delta_k=-y^2(\partial_u^2+\partial_y^2)+iky(\partial_u+i\partial_y).
\]
A function $h$ is holomorphic exactly when $\partial_{\bar\tau}h=0$, and
$\Delta_k=-\xi_{2-k}\xi_k$.

A \emph{harmonic Maass form} of weight $k$ is a real-analytic function on $\HH$
that transforms like a modular form of weight $k$ (with some multiplier),
satisfies $\Delta_kh=0$, and grows at most exponentially at every cusp. Such
an $h$ splits as a holomorphic part plus a non-holomorphic correction. The
holomorphic part is called a \emph{mock modular form}; the function $h$ is its
\emph{completion}; and $\xi_kh$, which is a holomorphic modular form of weight
$2-k$, is its \emph{shadow}. A product of a mock modular form and a modular
form (holomorphic on $\HH$, possibly with poles at the cusps) is called
\emph{mixed mock modular}; its weight is the sum of the two weights.

\begin{fact}[eta quotients {\cite[Proposition~2.5]{BS}}]\label{fact:etaquot}
Let $g(\tau)=\prod_{\delta\mid M}\eta(\delta\tau)^{r_\delta}$ with integers
$r_\delta$ whose sum $2k$ is even. If $\sum_\delta\delta r_\delta\equiv0$ and
$\sum_\delta(M/\delta)r_\delta\equiv0\pmod{24}$, then for every
$\gamma\in\Gamma_0(M)$,
$g(\gamma\tau)=\chi(d)(c\tau+d)^kg(\tau)$ with
$\chi(d)=\big(\frac{(-1)^ks}{d}\big)$, $s=\prod_\delta\delta^{r_\delta}$.
\end{fact}

The \emph{eta multiplier} $\nu_\eta(\gamma)$ is the root of unity defined by
$\eta(\gamma\sigma)=\nu_\eta(\gamma)(c\sigma+d)^{1/2}\eta(\sigma)$ for
$\gamma\in\mathrm{SL}_2(\Z)$, where the square root has argument in
$(-\pi/2,\pi/2]$.

\subsection{Zwegers' functions}\label{sec:zwegers}
For $\sigma\in\HH$ let $\Lambda_\sigma=\Z\sigma+\Z$. Following
\cite[\S1.3]{Zwegers}, for complex $v$ and $w$ put
\[
\vartheta(v;\sigma)=\sum_{\nu\in\Z+\frac12}e^{\pi i\nu^2\sigma+2\pi i\nu(v+\frac12)},
\qquad
\mu(w,v;\sigma)=\frac{e^{\pi iw}}{\vartheta(v;\sigma)}
\sum_{r\in\Z}\frac{(-1)^re^{\pi i(r^2+r)\sigma+2\pi irv}}{1-e^{2\pi i(r\sigma+w)}},
\]
the second for $w,v\notin\Lambda_\sigma$ (the zeros of $\vartheta(\cdot;\sigma)$
are exactly the points of $\Lambda_\sigma$). For real $t$ let
$E(t)=2\int_0^te^{-\pi s^2}\,ds$, and for $\sigma\in\HH$ let
\[
R(0;\sigma)=\sum_{\nu\in\Z+\frac12}\big(\sgn\nu-E(\nu\sqrt{2\Img\sigma})\big)
(-1)^{\nu-\frac12}e^{-\pi i\nu^2\sigma},
\]
where $\sgn t$ is $1$, $0$ or $-1$ as $t$ is positive, zero or negative. The
\emph{diagonal completion} of $\mu$ is
\[
\muh(w,w;\sigma)=\mu(w,w;\sigma)+\tfrac{i}{2}R(0;\sigma).
\]
The function $\mu$ is holomorphic but not modular; adding the correction
$\frac i2R$ restores modularity at the cost of holomorphy.

\begin{fact}[{\cite[Theorem~1.11]{Zwegers}}]\label{fact:zwegers}
For $\sigma\in\HH$ and $w\notin\Lambda_\sigma$:
\begin{enumerate}[label=(\alph*)]
\item $\muh(w+j\sigma+l,\,w+j\sigma+l;\sigma)=\muh(w,w;\sigma)$ for all integers
$j,l$;
\item $\muh\big(\tfrac{w}{c\sigma+d},\tfrac{w}{c\sigma+d};\gamma\sigma\big)
=\nu_\eta(\gamma)^{-3}(c\sigma+d)^{1/2}\muh(w,w;\sigma)$ for every
$\gamma\in\mathrm{SL}_2(\Z)$.
\end{enumerate}
\end{fact}

\section{The finite identity}\label{sec:finite}

In this section and the next, $q$ is real with $0<q<1$; \cref{thm:main} is
extended to complex $q$ at the end of \cref{sec:limit}. For an integer
$n\ge1$ write $x=q^{2n}$ and define
\begin{equation}\label{eq:f0}
f_0(n)=\frac{q\,B_n\,(1-x^4)^3(x^2-qx+q^2)}{2(1+x^4)(x-q)^2(q^2+x^4)},
\end{equation}
and, for each integer $k\ge0$,
\begin{equation}\label{eq:F}
F(n,k)=f_0(n)\Big(\frac qx\Big)^k\,\frac{1+q^3Q^{2k}/x^3}{1+q^3/x^3}\,
\frac{\poch{qx,\ Q/x^2,\ -Q/x^2,\ -Q/x^2}{Q}{k}}
{\poch{q^3/x,\ q^3/x,\ -q^3/x,\ -q^4/x^4}{Q}{k}} .
\end{equation}
No denominator vanishes: $1-q^3Q^j/x=1-q^{3+2j-2n}$ has an odd exponent, and
the other factors are of the form $1+(\text{positive})$. The factor
$\poch{Q/x^2}{Q}{k}=\prod_{j<k}(1-q^{2j+2-4n})$ contains the factor
$1-q^{0}=0$ (from $j=2n-1$) as soon as $k\ge2n$, so $F(n,k)=0$ for $k\ge2n$.
Hence
\[
T_n=\sum_{k=0}^{2n-1}F(n,k)=\sum_{k\ge0}F(n,k)
\]
is a finite sum.

\begin{lemma}\label{lem:telescope}
For every integer $N\ge1$,
\begin{equation}\label{eq:telescope}
\sum_{n=0}^{N-1}a_n=-T_N .
\end{equation}
\end{lemma}

\begin{proof}
Let $\cert(q,x,y)$ be the rational function defined in \cref{app:cert}, and
put
\[
G(n,k)=\cert(q,q^{2n},q^{2k})\,F(n,k).
\]
The identities proved in \cref{app:cert} say exactly that
\begin{equation}\label{eq:wz}
F(n+1,k)-F(n,k)=G(n,k+1)-G(n,k),\qquad G(n,0)=a_n,
\end{equation}
so $(F,G)$ is a Wilf--Zeilberger pair. Two values of $k$ need care. As a
function of $y=q^{2k}$, the certificate $\cert$ has denominator factors
$x^2-y$ and $q^2x^2-y$, which vanish at $k=2n$ and $k=2n+1$; at exactly these
$k$, the numerator of $F(n,k)$ contains the vanishing factor $1-Q^{2n}/x^2$.
To give $G(n,2n)$ and $G(n,2n+1)$ a meaning, treat $x$ as a real variable near
$q^{2n}$, so that both factors vanish to first order as $x\to q^{2n}$. (The
factor $B_n$ is continued in $x$ by
$\poch{a}{r}{z}=\poch{a}{r}{\infty}/\poch{ar^z}{r}{\infty}$, $r^z=e^{z\log r}$;
it stays finite and nonzero.) The ratio of the two vanishing factors has a
finite nonzero limit, and $G(n,k)$ is defined as the limit of the product.
Every other denominator factor of $\cert$ is nonzero for $0<q<1$. Since
\eqref{eq:wz} holds as an identity of analytic functions of $x$ away from these
points, it continues to hold in the limit.

For $k\ge2n+2$ the certificate is finite and $F(n,k)=0$, so $G(n,k)=0$.
Summing \eqref{eq:wz} over $0\le k\le2n+1$, the right side telescopes to
$G(n,2n+2)-G(n,0)=-a_n$, while the left side is $T_{n+1}-T_n$ (recall
$F(n,2n)=F(n,2n+1)=0$). The two values $G(n,2n)$ and $G(n,2n+1)$ cancel inside
the telescoping sum and are never set to zero. Thus $T_{n+1}-T_n=-a_n$ for
$n\ge1$. Finally, direct substitution in \eqref{eq:F} gives
$T_1=F(1,0)+F(1,1)=-(1-q^5)/[(1+q)(1+q^2)^2]=-a_0$ by \eqref{eq:a0}, and
adding the recurrences gives \eqref{eq:telescope}.
\end{proof}

The finite sum $T_N$ is a terminating ${}_8\phi_7$, so Watson's transformation
applies. For an integer $N\ge1$ put $x=q^{2N}$ and define
\begin{align}
C_N&=\frac{q^2\,x\,B_N\,(1-x^4)^3}{2(1+x^4)(x-q)^2(q^2+x^4)},\label{eq:CN}\\
h_N(k)&=C_N\,Q^k\,
\frac{\poch{Q/x^2}{Q}{k}^2\,\poch{-Q/x^2}{Q}{k}}
{\poch{-Q^2/x^4}{Q}{k}\,\poch{q^3/x}{Q}{k}^2}\qquad(0\le k\le2N-1).\label{eq:hN}
\end{align}

\begin{lemma}\label{lem:watson}
For every integer $N\ge1$,
\begin{equation}\label{eq:watsonsum}
\sum_{n=0}^{N-1}a_n=\sum_{k=0}^{2N-1}h_N(k),
\end{equation}
and every term $h_N(k)$ is nonnegative.
\end{lemma}

\begin{proof}
Apply \cref{fact:watson} with $m=2N-1$ and
\[
a=-q^3/x^3,\qquad b=qx,\qquad c=e=-Q/x^2,\qquad d=Q .
\]
Then $Q^{-m}=Q/x^2$. On the left of \cref{fact:watson}, the parameter $a$
cancels against $aQ/d=a$, and $d=Q$ cancels against the factor $\poch{Q}{Q}{k}$
in the definition of ${}_8\phi_7$. The pair $\poch{\pm Q\sqrt a}{Q}{k}/\poch{\pm\sqrt a}{Q}{k}$
equals $(1-aQ^{2k})/(1-a)=(1+q^3Q^{2k}/x^3)/(1+q^3/x^3)$, and the argument is
$a^2Q^{m+2}/(bcde)=q/x$. Comparing with \eqref{eq:F}, the left side is
$T_N/f_0(N)$. On the right, the prefactor is
\[
\begin{aligned}
\frac{\poch{-q^5/x^3,\ q/x}{Q}{m}}{\poch{-q^3/x^3,\ q^3/x}{Q}{m}}
&=\frac{1+q^3Q^m/x^3}{1+q^3/x^3}\cdot\frac{1-q/x}{1-qQ^m/x}
=\frac{1+q/x}{1+q^3/x^3}\cdot\frac{1-q/x}{1-x/q}\\
&=\frac{x^2}{x^2-qx+q^2}\Big(-\frac qx\Big)
=-\frac{qx}{x^2-qx+q^2},
\end{aligned}
\]
because each ratio telescopes to a single quotient, $Q^m=x^2/q^2$, and
$x^3+q^3=(x+q)(x^2-qx+q^2)$.
In the ${}_4\phi_3$, the parameter $d=Q$ again cancels $\poch{Q}{Q}{k}$. Writing
out the remaining parameters $aQ/(bc)=Q/x^2$, $e=-Q/x^2$, $Q^{-m}=Q/x^2$,
$aQ/b=-Q^2/x^4$, $aQ/c=q^3/x$ and $deQ^{-m}/a=q^3/x$,
we find that $-f_0(N)$ times the prefactor times the $k$-th term of the
${}_4\phi_3$ is exactly $h_N(k)$. Combining with \cref{lem:telescope} gives
\eqref{eq:watsonsum}.

For nonnegativity, note that in \eqref{eq:CN} and \eqref{eq:hN} the factors
that could be negative, $\poch{Q/x^2}{Q}{k}$ and $\poch{q^3/x}{Q}{k}$, occur
squared, and the latter is nonzero. Every other factor is positive when
$0<q<1$ (note $1-x^4>0$).
\end{proof}

\section{The centred limit and the proof of \texorpdfstring{\cref{thm:main}}{Theorem 1.1}}\label{sec:limit}

We keep $0<q<1$. The terms $h_N(k)$ are largest near the middle $k\approx N$ of
their range, so we re-index around it. For every integer $r$ put
\begin{equation}\label{eq:br}
b_r=\frac{q^{2r(r+1)}}{\poch{q}{Q}{r+1}^2},
\end{equation}
using the negative-index convention \eqref{eq:negindex} when $r+1<0$. Unwinding
that convention gives, for $j\ge0$,
\begin{equation}\label{eq:bneg}
b_{-j-1}=q^{2j}\,\poch{q}{Q}{j}^2 .
\end{equation}
For integers $N\ge1$ and $-N\le r\le N-1$ define
\begin{equation}\label{eq:kappa}
\kappa_N=\frac{C_N\,q^{2-2N}}{\poch{q}{Q}{N-1}^2},\qquad
\rho_{N,r}=\frac{\poch{q^{2N-2r}}{Q}{N+r}^2\,\poch{-q^{2N-2r}}{Q}{N+r}}
{\poch{-q^{6N-2r-2}}{Q}{N+r}} .
\end{equation}

\begin{lemma}\label{lem:centre}
For $N\ge1$ and $-N\le r\le N-1$,
\begin{equation}\label{eq:centre}
h_N(N+r)=\kappa_N\,b_r\,\rho_{N,r}.
\end{equation}
Moreover,
\begin{equation}\label{eq:Sb}
S(q)=\frac{\poch{q}{Q}{\infty}^2J_2}{2J_4^4}\sum_{r\in\Z}b_r .
\end{equation}
\end{lemma}

\begin{proof}
\emph{The identity \eqref{eq:centre}.} At $r=-N$ the left side is
$h_N(0)=C_N$. On the right, $\rho_{N,-N}=1$ (all its products are empty) and
$b_{-N}=q^{2N-2}\poch{q}{Q}{N-1}^2$ by \eqref{eq:bneg}, so the right side is
also $C_N$. It remains to compare the ratios of consecutive terms. Put
$x=q^{2N}$ and $v=q^{2r}$. From \eqref{eq:hN},
\[
\frac{h_N(N+r+1)}{h_N(N+r)}=\frac{Q(1-Qv/x)^2(1+Qv/x)}{(1+Q^2v/x^3)(1-q^3v)^2},
\]
and from \eqref{eq:br} and \eqref{eq:kappa},
\[
\frac{b_{r+1}\rho_{N,r+1}}{b_r\rho_{N,r}}
=\frac{Q^2v^2}{(1-q^3v)^2}\cdot\frac{(1-x/(Qv))^2(1+x/(Qv))}{1+x^3/(Q^2v)} .
\]
After the inner fractions are cleared, both ratios equal
\[
\frac{Q\,(x-Qv)^2(x+Qv)}{(x^3+Q^2v)(1-q^3v)^2},
\]
and induction on $r$ gives \eqref{eq:centre}.

\emph{Bounds.} For each fixed $r$, the first arguments $\pm q^{2N-2r}$ and
$-q^{6N-2r-2}$ in \eqref{eq:kappa} tend to $0$ as $N\to\infty$, so
$\rho_{N,r}\to1$. Uniformly in $-N\le r\le N-1$,
\[
0\le\rho_{N,r}\le\poch{-q^2}{Q}{\infty},
\]
because $2N-2r\ge2$, the factors $1-q^{2N-2r}Q^j$ lie in $(0,1]$, the factors
$1+q^{2N-2r}Q^j$ are at most $1+q^2Q^j$, and the denominator factors are at
least $1$. Next let $B_\infty=\lim_{N\to\infty}B_N$. Letting $N\to\infty$ in
\eqref{eq:Bn} and using $\poch{Q}{Q^2}{\infty}=J_2/J_4$ and
$\poch{-Q}{Q}{\infty}=J_4/J_2$,
\[
B_\infty=\frac{\poch{q}{Q}{\infty}^4J_2}{J_4^4},\qquad
\kappa_N\longrightarrow\frac{B_\infty}{2\poch{q}{Q}{\infty}^2};
\]
for the second limit note that $C_N\sim xB_N/(2q^2)$ as $x=q^{2N}\to0$. In
particular $(\kappa_N)_{N\ge1}$ is bounded. Finally, $b_r\le
q^{2r(r+1)}/\poch{q}{Q}{\infty}^2$ for $r\ge0$, and $b_{-j-1}\le q^{2j}$ for
$j\ge0$ by \eqref{eq:bneg}; so $\sum_r b_r$ converges absolutely.

\emph{The limit.} Extend $r\mapsto h_N(N+r)$ by zero outside
$-N\le r\le N-1$. The bounds just proved give
$0\le h_N(N+r)\le(\sup_M\kappa_M)\,\poch{-q^2}{Q}{\infty}\,b_r$, a summable
majorant that does not depend on $N$. By \cref{lem:watson} and dominated
convergence,
\[
S(q)=\lim_{N\to\infty}\sum_{r}h_N(N+r)=\frac{B_\infty}{2\poch{q}{Q}{\infty}^2}\sum_{r\in\Z}b_r,
\]
which is \eqref{eq:Sb}.
\end{proof}

\begin{proof}[Proof of \cref{thm:main}]
\emph{Real $q$.} Let $0<q<1$. By \eqref{eq:bneg}, the first sum in
\cref{fact:LO} is $\sum_{r\le-2}b_r=\sum_{r<0}b_r-1$, and after the shift
$j=r+1$ the second sum is $1+\sum_{r\ge0}b_r$. Hence \cref{fact:LO} says
\begin{equation}\label{eq:bU}
\sum_{r\in\Z}b_r=-\frac{U(q)}{qJ_2}.
\end{equation}
Next split $U(q)$ by the parity of $r$. The odd terms $r=2s+1$ give
$-q^{4s^2+4s+1}/(1-q^{4s+1})$. In the even terms $r=2s$ replace $s$ by $-s$;
since $1/(1-q^{-4s-1})=-q^{4s+1}/(1-q^{4s+1})$, they give the same expression.
Therefore
\begin{equation}\label{eq:UL}
U(q)=-2q\sum_{s\in\Z}\frac{q^{4s(s+1)}}{1-q^{4s+1}}=-2qL(q).
\end{equation}
All rearrangements are allowed because the series converge absolutely.
Substituting \eqref{eq:bU} and \eqref{eq:UL} in \eqref{eq:Sb} gives
$S(q)=\poch{q}{Q}{\infty}^2L(q)/J_4^4$, which is \eqref{eq:main}.

\emph{Complex $q$.} Fix $0<t<1$. On the disc $|q|\le t$, each finite product in
$B_n$ is bounded in absolute value by a constant independent of $n$ (for
example $|\poch{q}{Q}{n}|\le\poch{-t}{t^2}{\infty}$), and each denominator
factor is bounded away from zero (for example
$|\poch{Q^2}{Q^2}{2n}|\ge\poch{t^4}{t^4}{\infty}$). The same holds for the extra
factors $1+q^{8n}$, $1+q^{2n+1}$, $1+q^{4n+2}$ in $a_n$, and $|P(q,q^{2n})|$ is
bounded. Hence $a_n=O_t(t^{2n})$, and $S$ converges absolutely and locally
uniformly on the unit disc. For the right side, rewrite the terms with $r<0$
of $L$ as in \eqref{eq:UL}:
\[
L(q)=\sum_{r\ge0}\frac{q^{4r(r+1)}}{1-q^{4r+1}}-\sum_{r\ge1}\frac{q^{4r^2-1}}{1-q^{4r-1}} .
\]
This form converges locally uniformly on the unit disc and shows that $L$ is
holomorphic there, including at $q=0$, with $L(0)=1$; the products
$\poch{q}{q^2}{\infty}$ and $J_4$ are holomorphic and nonzero there. Both sides of
\eqref{eq:main} are therefore holomorphic on the unit disc and agree on
$(0,1)$, so they agree everywhere by the identity theorem. Finally $S(0)=a_0(0)=1$
by \eqref{eq:a0}.
\end{proof}

\subsection*{Relation with Au's Example~XXI}
Let $\SXXI(q)$ be the series on the left of \cite[(3.13)]{Au}, whose summand is
\[
\frac{q^{2n}\poch{q}{Q}{n}^6}{(1+q^{2n+1})^3\poch{Q}{Q}{2n}^3}
\big(3+3q^{2n+1}-q^{4n}+2q^{4n+2}-3q^{6n+1}-3q^{8n+2}-q^{10n+3}\big).
\]
Au evaluated $\SXXI$ in \cite[(3.13)]{Au}; rewritten with \eqref{eq:UL} in the
present notation, that evaluation reads $\SXXI(q)=2J_1^4L(q)/J_2^8$. (We have
also confirmed this identity directly, as a power series modulo $q^{220}$.) Since $\poch{q}{Q}{\infty}=J_1/J_2$, comparing with
\cref{thm:main} gives the following.

\begin{corollary}\label{cor:XXI}
For $0<|q|<1$,
\[
S(q)=\frac{J_2^6}{2J_1^2J_4^4}\,\SXXI(q).
\]
\end{corollary}

Thus the Watson series and Example~XXI share the same mock modular part
$L(q)$; they differ by an infinite product, that is, by a modular form.

\section{The classical limit}\label{sec:classical}

This section proves \cref{prop:classical} directly from the definition of
$S$, without using \cref{thm:main}. Throughout, $0<q<1$.

\begin{proof}[Proof of \cref{prop:classical}]
\emph{Termwise limits.} Write $P_n(q)=P(q,q^{2n})$. Expanding \eqref{eq:P} into
monomials and setting $x=q^{2n}$ gives
\begin{equation}\label{eq:Pn}
P_n(q)=\sum_j\alpha_j\,q^{s_j+t_jn},
\end{equation}
with integer coefficients $\alpha_j$ and nonnegative integers $s_j,t_j$ ($t_j$
is twice the power of $x$). Direct inspection of \eqref{eq:P} gives
\[
\sum_j\alpha_j=0,\qquad \sum_j\alpha_js_j=-20,\qquad\sum_j\alpha_jt_j=-168,
\]
so $P_n(1)=0$ and $-P_n'(1)=168n+20=4(42n+5)$. Hence
$P_n(q)/(1-q)\to4(42n+5)$ as $q\to1$. Each factor $1-q^m$ in $B_n$ behaves like
$m(1-q)$, and the numerator and denominator of $B_n$ contain $4n+2n$ and $6n$
such factors respectively, while $\poch{-Q}{Q}{4n}\to2^{4n}$. Using
$(2n)!=4^nn!\,(1/2)_n$ one obtains
\[
\lim_{q\to1}B_n=\frac{(1/2)_n^3}{(n!)^3\,64^n},\qquad
\lim_{q\to1}\frac{8a_n}{1-q}=\frac{8\cdot4(42n+5)}{2\cdot2\cdot2^3}\,
\frac{(1/2)_n^3}{(n!)^3\,64^n}=\frac{(1/2)_n^3}{(n!)^3\,64^n}(42n+5).
\]

\emph{A uniform bound.} We must justify summing these limits. Put
$t=q^{2n}$. Cancelling common factors in \eqref{eq:Bn} gives
\[
\frac{B_{n+1}}{B_n}=\Big(\frac{1-qt}{1-q^2t}\Big)^3
\frac{\prod_{j=1}^4(1+q^{2j}t^4)}{(1+qt)(1+q^2t)^3(1+q^2t^2)^3(1+q^4t^2)^3}
\le(1+q^{4n+4})^{-6}.
\]
For the inequality: the first factor is at most $1$; the four numerator
factors are at most $(1+qt)(1+q^2t)^3$ (because $q^2t^4\le qt$ and
$q^{2j}t^4\le q^2t$); and each of the two remaining groups in the denominator
is at least $(1+q^4t^2)^3=(1+q^{4n+4})^3$.

Write $q=e^{-h}$ with $0<h<1$, let $m=\lfloor1/h\rfloor$ be the largest
integer at most $1/h$, and let $\lambda=(1+e^{-4})^{-6}<1$. For $n\le m-1$ we
have $(4n+4)h\le4$, hence $q^{4n+4}\ge e^{-4}$ and $B_{n+1}/B_n\le\lambda$.
For larger $n$ the ratio is at most $1$. Since $B_0=1$,
\[
B_n\le\lambda^n\quad(0\le n\le m),\qquad B_n\le\lambda^m\quad(n>m).
\]
Next, $\sum_j\alpha_j=0$ lets us write $P_n(q)=-\sum_j\alpha_j(1-q^{s_j+t_jn})$,
and $1-q^e\le e(1-q)$ for integers $e\ge0$. From \eqref{eq:P} one computes
$\sum_j|\alpha_j|s_j=112$ and $\sum_j|\alpha_j|t_j=416$, so
\[
\frac{|P_n(q)|}{1-q}\le112+416n .
\]
The extra denominator factors of $a_n$ are at least $1$ and $q^{2n}\le1$.
Therefore $|8a_n/(1-q)|\le8(112+416n)\lambda^n$ for $n\le m$, while
\[
\sum_{n>m}\Big|\frac{8a_n}{1-q}\Big|\le8\lambda^m\sum_{n>m}(112+416n)e^{-2hn}
=O\big(h^{-2}\lambda^{\lfloor1/h\rfloor}\big),
\]
which tends to $0$ as $h\to0$.

\emph{Conclusion.} On the range $n\le m$ the terms are dominated by the summable
sequence $8(112+416n)\lambda^n$, which does not depend on $h$; so by dominated
convergence that part of the sum tends to
$\sum_n(1/2)_n^3(42n+5)/((n!)^364^n)$. The remaining part tends to $0$. By
\eqref{eq:classical} the limit is $16/\pi$.
\end{proof}

\begin{remark}
Numerically, $8S(e^{-h})/(1-e^{-h})-16/\pi$ is about $2.5h$ for small $h$; for
example it equals $2.5\times10^{-4}$ (to two digits) at $h=10^{-4}$.
\end{remark}

\section{The completion and its shadow}\label{sec:completion}

Throughout this section $\tau=u+iy\in\HH$ and $q=e^{2\pi i\tau}$; the notation
is that of \cref{sec:modular,sec:zwegers}. Define the eta quotient
\[
A(\tau)=\frac{\eta(\tau)^4}{\eta(2\tau)^3\,\eta(4\tau)^4}
=q^{-3/4}\frac{J_1^4}{J_2^3J_4^4},
\]
and the functions
\[
f(\tau)=-\tfrac{i}{2}\,\mu(\tau,\tau;2\tau),\qquad
\fh(\tau)=-\tfrac{i}{2}\,\muh(\tau,\tau;2\tau)=f(\tau)+\tfrac14R_0(\tau),\qquad
R_0(\tau)=R(0;2\tau).
\]
They have no poles, because $\tau\notin\Lambda_{2\tau}=2\Z\tau+\Z$. Finally put
\[
\Sh(\tau)=A(\tau)\fh(\tau)=A(\tau)f(\tau)+\tfrac14A(\tau)R_0(\tau).
\]
In this language \cref{thm:completion} states: $S(e^{2\pi i\tau})=Af$; the
function $\Sh$ transforms like a modular form of weight $-1$ and character
$\chi_{-4}$ on $\Gamma_0(4)$; and $\xi_{1/2}\fh=-\frac12\eta(2\tau)^3$. In
the decomposition $S=A\cdot f$, the factor $A$ is a modular form of weight
$-3/2$ and $f$ is a mock modular form of weight $1/2$ with completion $\fh$.

\begin{proof}[Proof of \cref{thm:completion}]
\emph{The identity $S=Af$.} Substituting $v=w=\tau$ and $\sigma=2\tau$ in the
definitions of \cref{sec:zwegers} and re-indexing the sums ($\nu=m+\frac12$ in
$\vartheta$) gives
\[
\vartheta(\tau;2\tau)=-iq^{-1/4}\Theta(q),\qquad
\mu(\tau,\tau;2\tau)=-\frac{iq^{-1/4}}{\Theta(q)}\,U(q)=\frac{2iq^{3/4}}{\Theta(q)}\,L(q),
\]
the last step by \eqref{eq:UL}. Hence $f=q^{3/4}L(q)/\Theta(q)$. With
$\Theta=J_1^2/J_2$ (\cref{fact:jacobi}) and $\poch{q}{q^2}{\infty}=J_1/J_2$,
\[
A f=\frac{J_1^4}{J_2^3J_4^4}\cdot\frac{J_2}{J_1^2}\,L(q)
=\frac{\poch{q}{q^2}{\infty}^2}{J_4^4}\,L(q)=S(q)
\]
by \cref{thm:main}.

\emph{The shadow.} Replace $2\nu$ by the odd integer $\ell$ in the definition of
$R(0;2\tau)$:
\[
R_0(\tau)=\sum_{\ell\in2\Z+1}(-1)^{(\ell-1)/2}\big(\sgn\ell-E(\ell\sqrt y)\big)q^{-\ell^2/4}.
\]
The terms decay like a Gaussian in $\ell$ uniformly on compact subsets of
$\HH$, so the series may be differentiated term by term. Only $E(\ell\sqrt y)$
depends on $y$, and $\partial_yE(\ell\sqrt y)=\ell y^{-1/2}e^{-\pi\ell^2y}$;
the holomorphic factor $q^{-\ell^2/4}$ is killed by $\partial_{\bar\tau}$. Since
$e^{-\pi\ell^2y}q^{-\ell^2/4}=\overline{q^{\ell^2/4}}$, this gives
\[
\partial_{\bar\tau}R_0=-\frac{i}{2\sqrt y}\,
\overline{\sum_{\ell\in2\Z+1}(-1)^{(\ell-1)/2}\ell\,q^{\ell^2/4}} .
\]
The terms with $\ell$ and $-\ell$ are equal, so by \cref{fact:jacobi} the sum
is $2q^{1/4}J_2^3=2\eta(2\tau)^3$. Therefore
\begin{equation}\label{eq:dR}
\partial_{\bar\tau}R_0=-\frac{i}{\sqrt y}\,\overline{\eta(2\tau)^3},\qquad
\partial_{\bar\tau}\fh=\tfrac14\partial_{\bar\tau}R_0,\qquad
\partial_{\bar\tau}\Sh=-\frac{iA(\tau)}{4\sqrt y}\,\overline{\eta(2\tau)^3},
\end{equation}
using that $f$ and $A$ are holomorphic. Applying $\xi_{1/2}=2iy^{1/2}\overline{\partial_{\bar\tau}(\cdot)}$
to the middle identity gives $\xi_{1/2}\fh=-\frac12\eta(2\tau)^3$. Since this is
holomorphic, $\Delta_{1/2}\fh=-\xi_{3/2}\xi_{1/2}\fh=0$, so $\fh$ is harmonic.

\emph{The transformation law.} Let $\gamma=\left(\begin{smallmatrix}a&b\\c&d\end{smallmatrix}\right)\in\Gamma_0(2)$
and put $\gamma'=\left(\begin{smallmatrix}a&2b\\c/2&d\end{smallmatrix}\right)\in\mathrm{SL}_2(\Z)$,
so that $\gamma'(2\tau)=2\gamma\tau$ and $\frac c2(2\tau)+d=c\tau+d$. Apply
\cref{fact:zwegers}(b) with $\sigma=2\tau$, the matrix $\gamma'$, and
$w=a\tau+b$; then $w/(c\tau+d)=\gamma\tau$ and
\[
\muh(\gamma\tau,\gamma\tau;2\gamma\tau)=\nu_\eta(\gamma')^{-3}(c\tau+d)^{1/2}
\muh(a\tau+b,a\tau+b;2\tau).
\]
Since $a$ is odd, $a\tau+b=\tau+\frac{a-1}{2}(2\tau)+b$ differs from $\tau$ by an
element of $\Lambda_{2\tau}$, so by \cref{fact:zwegers}(a) the last factor is
$\muh(\tau,\tau;2\tau)$. On the other hand
$\eta(2\gamma\tau)^3=\eta(\gamma'(2\tau))^3=\nu_\eta(\gamma')^3(c\tau+d)^{3/2}\eta(2\tau)^3$.
Hence
\[
K(\tau)=\eta(2\tau)^3\,\muh(\tau,\tau;2\tau)
\]
satisfies $K(\gamma\tau)=(c\tau+d)^2K(\tau)$ for all $\gamma\in\Gamma_0(2)$:
it has weight $2$ and trivial character on $\Gamma_0(2)$.

Now let
\[
g(\tau)=\frac{\eta(\tau)^4}{\eta(2\tau)^6\,\eta(4\tau)^4},
\]
so that $A=g\,\eta(2\tau)^3$ and $\Sh=-\frac i2\,g\,K$. The exponents
$r_1=4$, $r_2=-6$, $r_4=-4$ satisfy $\sum_\delta\delta r_\delta=-24$ and
$\sum_\delta(4/\delta)r_\delta=0$, and $\frac12\sum_\delta r_\delta=-3$; also
$s=\prod\delta^{r_\delta}=2^{-14}$, so $\big(\frac{(-1)^{-3}s}{d}\big)=\big(\frac{-1}{d}\big)=\chi_{-4}(d)$.
By \cref{fact:etaquot}, $g$ has weight $-3$ and character $\chi_{-4}$ on
$\Gamma_0(4)$. (If one prefers to apply the criterion only to eta products with
nonnegative exponents, write $g=\eta(\tau)^4\eta(2\tau)^{18}\eta(4\tau)^{20}/(\eta(2\tau)^{24}\eta(4\tau)^{24})$:
the numerator has weight $21$ and $s=2^{58}$, so character $\chi_{-4}$, and the
denominator $\Delta(2\tau)\Delta(4\tau)$ has weight $24$ and trivial character.)
Multiplying, $\Sh=-\frac i2gK$ has weight $-3+2=-1$ and character $\chi_{-4}$
on $\Gamma_0(4)\subset\Gamma_0(2)$. This is the transformation law of
\cref{thm:completion}. The same computation shows that $A=g\,\eta(2\tau)^3$ is
a modular form of weight $-3/2$ (with the product of the two multipliers), and
that $\fh=-\frac i2K/\eta(2\tau)^3$ transforms with weight $2-\frac32=\frac12$.

\emph{Growth at the cusps.} The group $\Gamma_0(4)$ has three cusps, $\infty$,
$0$ and $\frac12$; the cusp $\frac12$ is the image of $\infty$ under
$\left(\begin{smallmatrix}1&0\\2&1\end{smallmatrix}\right)\in\Gamma_0(2)$, so for
the $\Gamma_0(2)$-covariant function $\fh$ it suffices to look at $\infty$ and
$0$. At $\infty$, $f=q^{3/4}(1+O(q))$ because $L(0)=\Theta(0)=1$, and the Gaussian
decay of $1-|E|$ gives $R_0(\tau)=O(y^{-1/2}e^{-\pi y/2})$. For the cusp $0$, put
$\tau=-1/(2z)$ with $z\in\HH$, $v=\Img z$ and $p=e^{2\pi iz}$. Zwegers' inversion
law (\cref{fact:zwegers}(b) with $\gamma=\left(\begin{smallmatrix}0&-1\\1&0\end{smallmatrix}\right)$)
(where $\nu_\eta=e^{-\pi i/4}$, $w=-\frac12$), followed by \cref{fact:zwegers}(a)
with $j=0$, $l=1$, gives
\[
\muh\Big(-\frac1{2z},-\frac1{2z};-\frac1z\Big)=e^{3\pi i/4}z^{1/2}\,\muh\big(\tfrac12,\tfrac12;z\big).
\]
From the definitions, $\vartheta(\frac12;z)=-2p^{1/8}(1+O(p))$ (the terms
$\nu=\pm\frac12$ dominate), and
\[
\sum_{r\in\Z}\frac{p^{r(r+1)/2}}{1+p^r}=\frac12+O(p),
\]
since the term $r=0$ is $\frac12$ and the terms $r$ and $-r$ combine to
$2p^{r(r+1)/2}/(1+p^r)$ for $r\ge1$. Hence
$\mu(\frac12,\frac12;z)=-\frac i4p^{-1/8}(1+O(p))$, while
$R(0;z)=O(v^{-1/2}e^{-\pi v/4})$. So $\fh$ grows at most exponentially at both
cusps of $\Gamma_0(2)$. An eta quotient has finite exponential order at every
cusp, so the same holds for $A$ and for $\Sh$. This completes the proof.
\end{proof}

\begin{remark}
The shadow belongs to the weight-$1/2$ factor $\fh$, not to $\Sh$ itself. By
\eqref{eq:dR},
\[
\xi_{-1}\Sh=-\frac{\overline{A(\tau)}}{2y^{3/2}}\,\eta(2\tau)^3,
\]
which is not holomorphic. So $\Sh$ is \emph{not} a harmonic Maass form of
weight $-1$; it is the product of the harmonic Maass form $\fh$ with the
modular form $A$, which is what ``mixed mock modular'' means.
\end{remark}

\section{Comparison with Campbell's \texorpdfstring{$q$}{q}-analogue}\label{sec:campbell}

Campbell \cite[Theorem~2]{Campbell1} gives another $q$-analogue of
\eqref{eq:classical}. We show that its summand is an elementary multiple of
the summand of Au's Example~III, which is a modular (not mock modular)
evaluation. For each integer $n\ge0$ put
\[
w_n=1+3q^{2n+1}+3q^{4n+2}-2q^{6n+1}+q^{6n+3}-3q^{8n+2}-3q^{10n+3},
\qquad
u_n=\frac{q^{6n^2}\,w_n\,\poch{q}{q^2}{n}^6}{(1+q^{2n+1})^3\,\poch{q^2}{q^2}{2n}^3}.
\]
The term $u_n$ is the summand of Au's Example~III. In this section only, let
$c_n$ denote Campbell's summand,
\[
c_n=\frac{(-1)^nq^{(n-1)(3n+5)}\poch{q}{q^2}{n}^4\poch{q^3}{q^2}{n}}
{\poch{-q^2}{q^2}{n}^2\poch{q^2}{q^2}{n}^3\poch{q^{3-4n}}{q^2}{n}\poch{q^{2n+2}}{q^2}{n}}
\cdot\frac{\poch{q^{-1},q}{q^4}{n}}{\poch{q^2,q^6}{q^4}{n}}
\cdot\frac{-(1-q^{2n-1})\,w_n}{(1+q^{2n+1})^2}.
\]

\begin{proposition}\label{prop:campbell}
For $0<|q|<1$ and every integer $n\ge0$,
\[
c_n=\frac{1-q^2}{q^6}\,u_n .
\]
\end{proposition}

\begin{proof}
Reversing the order of the factors in the product with negative exponents,
\[
\poch{q^{3-4n}}{q^2}{n}=(-1)^nq^{-3n^2+2n}\poch{q^{2n-1}}{q^2}{n}.
\]
Separating even and odd factors (\cref{fact:evenodd}),
\[
\poch{q^{-1}}{q^4}{n}\poch{q}{q^4}{n}=\poch{q^{-1}}{q^2}{2n},\qquad
\frac{(1-q^{2n-1})\poch{q^{-1}}{q^2}{2n}}{\poch{q^{2n-1}}{q^2}{n}}=(1-q^{-1})\poch{q}{q^2}{n}.
\]
The remaining cancellations use
\[
\poch{-q^2}{q^2}{n}\poch{q^2}{q^2}{n}=\poch{q^4}{q^4}{n},\qquad
\poch{q^2}{q^4}{n}\poch{q^4}{q^4}{n}=\poch{q^2}{q^2}{2n},
\]
\[
\frac{\poch{q^3}{q^2}{n}}{\poch{q}{q^2}{n}}=\frac{1-q^{2n+1}}{1-q},\qquad
\frac{\poch{q^6}{q^4}{n}}{\poch{q^2}{q^4}{n}}=\frac{1-q^{4n+2}}{1-q^2}.
\]
Substituting all of these into $c_n$ leaves $(1-q^2)u_n/q^6$. For $n=0$ all
products are empty and the identity is checked directly.
\end{proof}

Let $F_C$ and $G_C$ be the Wilf--Zeilberger pair used by Campbell
\cite[(41)--(43)]{Campbell1}. He obtains $c_n$ by forming
$F_C(n+1,n+k)+G_C(n,n+k)$ and replacing $q$ by $q^2$. So
\cref{prop:campbell} identifies the transformed summand itself, not only its
classical limit: Campbell's series is, up to the factor $(1-q^2)/q^6$, Au's
modular Example~III, and not the mock modular Watson series studied here.

Campbell's later multiparameter construction \cite[\S4.4, (32)]{Campbell2}
uses, with its base renamed $Q$, the kernel
\[
M_C(n,k)=Q^k\frac{\poch{Q^{n+\alpha},Q^{n+\beta}}{Q}{k}}{\poch{Q^{2n+\gamma},Q^{2n+\delta}}{Q}{k}},
\]
with complex parameters $\alpha,\beta,\gamma,\delta$ (for $0<q<1$ interpret
$Q^\alpha=e^{\alpha\log Q}$). Its ratio $M_C(n,k+1)/M_C(n,k)$ is a rational
function of $Q^k$ whose numerator and denominator have degree at most $2$. For
our kernel \eqref{eq:F}, the ratio $F(n,k+1)/F(n,k)$ is the function
$r^{(k)}$ of \cref{app:cert}; in lowest terms its numerator and denominator
both have degree $6$ in $Q^k$. Multiplying a kernel by a factor depending only
on $n$ does not change this ratio, and specializing the parameters cannot raise
its degree from $2$ to $6$, so the two kernels are not related in this simple way (with the same
base and index). This does not exclude more elaborate transformations
between the two sums.

\appendix
\section{The Wilf--Zeilberger certificate}\label{app:cert}

In this appendix $q$, $x$ and $y$ are independent variables; in
\cref{sec:finite} they are specialized to $q$, $x=q^{2n}$ and $y=q^{2k}=Q^k$.
Define the polynomial
\[
\begin{aligned}
D(q,x,y)={}&(x-1)^2(x+1)^2(x^2+1)(x^2-y)(x^2+y)^2\\
&\times(qx+1)(q^2x^2+1)(q^2x^2-y)(q^2x^2+y)^2,
\end{aligned}
\]
the rational function
\[
\Phi(q,x,y)=\frac{(x-qy)^2(x+qy)(q^2y+x^4)}{q\,y\,D(q,x,y)\,(q^3y^2+x^3)},
\]
and the polynomial $V(q,x,y)=\sum_{j=0}^5v_j(q,x)y^j$ with
\[
\begin{aligned}
v_0={}&q^9x^{15},\\
v_1={}&q^9x^{13}+q^8x^{14}-q^8x^{10}+q^7x^{13}+q^7x^{11}-q^7x^9+q^6x^{12}-q^6x^8\\
&+q^5x^{11}-q^5x^9-q^4x^{10},\\
v_2={}&q^8x^{10}-q^8x^8-q^7x^{13}+q^7x^{11}+q^6x^{10}-q^6x^6+q^5x^9-q^5x^5\\
&+q^4x^8-2q^4x^6-q^2x^6,\\
v_3={}&-q^7x^9-2q^5x^9+q^5x^7-q^4x^{10}+q^4x^6-q^3x^9+q^3x^5+q^2x^4\\
&-q^2x^2-qx^7+qx^5,\\
v_4={}&-q^5x^5-q^4x^6+q^4x^4-q^3x^7+q^3x^3-q^2x^6+q^2x^4+q^2x^2\\
&-qx^5+qx+x^2,\\
v_5={}&1.
\end{aligned}
\]
The certificate is
\[
\cert(q,x,y)=x\,(y+1)\,V(q,x,y)\,\Phi(q,x,y).
\]
Define also
\[
\mathcal N(q,x,y)=(x^4+y)(qx-1)(q^2x^4+y)(q^3x^3+y^2)(q^4x^4+y)(qxy-1),
\]
and the two ratio functions
\[
r^{(n)}(q,x,y)=q^{-1}\mathcal N(q,x,y)\,\Phi(q,x,y),
\]
\[
r^{(k)}(q,x,y)=\frac qx\cdot\frac{1+q^7y^2/x^3}{1+q^3y^2/x^3}\cdot
\frac{(1-qxy)(1-q^2y/x^2)(1+q^2y/x^2)^2}{(1-q^3y/x)^2(1+q^3y/x)(1+q^4y/x^4)} .
\]
Clearing the inner fractions,
\[
r^{(k)}=\frac{q\,(x^3+q^7y^2)(1-qxy)(x^2-q^2y)(x^2+q^2y)^2}
{(x^3+q^3y^2)(x-q^3y)^2(x+q^3y)(x^4+q^4y)},
\]
whose numerator and denominator have degree $6$ in $y$ and no common factor.

\emph{Meaning of the ratios.} Cancelling the finite products in \eqref{eq:F}
shows that $r^{(k)}=F(n,k+1)/F(n,k)$ after the substitution $x=q^{2n}$,
$y=q^{2k}$. Similarly $r^{(n)}=F(n+1,k)/F(n,k)$. One way to see the latter is to
check it at $y=1$, where it reduces to $f_0(n+1)/f_0(n)$, computed from
\eqref{eq:f0} and the ratio $B_{n+1}/B_n$ of the products in \eqref{eq:Bn},
and then use the compatibility identity
\begin{equation}\label{eq:compat}
r^{(n)}(q,x,y)\,r^{(k)}(q,q^2x,y)=r^{(k)}(q,x,y)\,r^{(n)}(q,x,q^2y),
\end{equation}
which expresses that the two ways of going from $F(n,k)$ to $F(n+1,k+1)$ agree.

\emph{The certificate identity.} The identity
\begin{equation}\label{eq:A1}
r^{(k)}(q,x,y)\,\cert(q,x,q^2y)-\cert(q,x,y)=r^{(n)}(q,x,y)-1
\end{equation}
holds in the field of rational functions of $q,x,y$. Multiplying by the common
denominator turns it into a polynomial identity, which is verified by
expansion (after clearing denominators it has $7278$ monomials; see the
verification scripts). Multiplying \eqref{eq:A1} by $F(n,k)$ and specializing
gives the first identity in \eqref{eq:wz}.

\emph{The boundary value.} The identity
\[
\cert(q,x,1)\cdot\frac{q(1-x^4)^3(x^2-qx+q^2)}{2(1+x^4)(x-q)^2(q^2+x^4)}
=\frac{x\,P(q,x)}{(1+x^4)(1+qx)(1+q^2x^2)^3}
\]
is again verified by polynomial expansion. Multiplied by $B_n$ and specialized
to $x=q^{2n}$ it says $G(n,0)=\cert(q,q^{2n},1)F(n,0)=a_n$, the second identity
in \eqref{eq:wz}. The removable singularities that occur in the
specialization were treated in the proof of \cref{lem:telescope}.


\section*{Declarations}
\noindent\textbf{Funding.} No external funding was received for this work.

\noindent\textbf{Competing interests.} The author declares no competing interests.

\noindent\textbf{Data availability.} The results are established by the proofs
in this article; no external data sets are used. Scripts that check every
identity in the paper (exactly where possible) are available at
\url{https://github.com/ramanujamath/au-watson-mock-series}.

\begin{thebibliography}{9}

\bibitem{Au}
K.~C.~Au, \emph{$q$-analogues of Wilf--Zeilberger seeds and Ramanujan
$1/\pi^k$-formulas}, Adv.\ Math.\ \textbf{481} (2026), 111181,
doi:10.1016/j.aim.2026.111181. Equation and example numbers follow
arXiv:2403.04555v4.

\bibitem{Campbell1}
J.~M.~Campbell, \emph{$q$-analogues of $\pi$-formulas due to Ramanujan and
Guillera}, J.\ Symbolic Comput.\ \textbf{134} (2026), 102531,
doi:10.1016/j.jsc.2025.102531. Numbering follows arXiv:2509.05897v1.

\bibitem{Campbell2}
J.~M.~Campbell, \emph{The construction of $q$-analogues via ${}_3\phi_2$-series
and $q$-difference equations}, arXiv:2602.14314v1, 2026.

\bibitem{DLMF}
\emph{NIST Digital Library of Mathematical Functions}, \S17.9, equation
(17.9.15), \url{https://dlmf.nist.gov/17.9}.

\bibitem{LO}
J.~Lovejoy and R.~Osburn, \emph{Mixed mock modular $q$-series}, J.\ Indian
Math.\ Soc.\ (N.S.), Special volume commemorating the 125th birth anniversary
of Srinivasa Ramanujan (2013), 45--61; arXiv:1304.1905.

\bibitem{Zwegers}
S.~Zwegers, \emph{Mock theta functions}, Ph.D.\ thesis, Universiteit Utrecht,
2002; arXiv:0807.4834.

\bibitem{BS}
M.~Boylan and Swati, \emph{Congruence properties modulo prime powers for a
class of partition functions}, Res.\ Number Theory \textbf{12} (2026), article
29, doi:10.1007/s40993-026-00715-4.

\end{thebibliography}

\end{document}
